\documentclass[11pt,a4paper]{article}
\usepackage[utf8]{inputenc}
\usepackage[T1]{fontenc}
\usepackage{lmodern}
\usepackage[french]{babel}
\usepackage{amsmath,amssymb}
\usepackage{geometry}
\geometry{top=1.7cm,bottom=1.7cm,left=1.7cm,right=1.7cm}
\usepackage{tikz}
\usetikzlibrary{arrows.meta}
\usepackage{enumitem}
\usepackage{xcolor}
\usepackage{array}
\usepackage{fancyhdr}
\usepackage[most]{tcolorbox}

\definecolor{bleu}{RGB}{0,84,150}
\definecolor{orange}{RGB}{225,110,0}
\definecolor{vert}{RGB}{0,125,85}

\newtcolorbox{encadre}[2][]{colback=#2!6,colframe=#2,boxrule=0.8pt,arc=2pt,
  left=6pt,right=6pt,top=5pt,bottom=5pt,fonttitle=\bfseries,title=#1}

\pagestyle{fancy}
\fancyhf{}
\renewcommand{\headrulewidth}{0.4pt}
\fancyhead[L]{\small\color{bleu}\textbf{1\textsuperscript{re} STI2D -- Nombres complexes}}
\fancyhead[R]{\small\color{bleu}Activité de découverte}
\fancyfoot[C]{\small\thepage}

\setlist[enumerate]{itemsep=3pt,topsep=3pt}
\setlist[enumerate,2]{itemsep=1pt,topsep=2pt}


\usepackage{multicol}

\begin{document}

\begin{center}
{\Large\textbf{Activité de découverte}}\\[3pt]
{\LARGE\textbf{Les nombres complexes}}\\[5pt]
{\large Première STI2D \quad$\bullet$\quad Durée : 1\,h \quad$\bullet$\quad Calculatrice autorisée}\\[3pt]
{\normalsize\itshape Forme algébrique (F.\,alg.) \quad et \quad forme trigonométrique (F.\,tr.)}
\end{center}

\begin{encadre}[Objectifs et compétences]{bleu}
\begin{multicols}{2}
\textbf{Objectifs :}
\begin{itemize}[leftmargin=1.2em,itemsep=1pt,topsep=1pt]
  \item découvrir un nouvel ensemble de nombres ;
  \item passer d'un point du plan à un nombre et inversement ;
  \item découvrir module, argument et forme trigonométrique ;
  \item appliquer ces outils à un contexte STI2D.
\end{itemize}
\columnbreak

\textbf{Compétences :} Chercher -- Représenter -- Calculer -- Communiquer.\\[4pt]
\textbf{Prérequis :} repérage orthonormé, théorème de Pythagore, cosinus, sinus et tangente.
\end{multicols}
\end{encadre}

%=====================================================================
\section*{Partie A -- Du point au nombre : la forme algébrique (F.\,alg.)}

Dans tout le repère orthonormé $(O;\vec{u},\vec{v})$ ci-dessous, à tout point $M(x\,;y)$
on associe le nombre $z_M = x + iy$, où $i$ est un nombre tel que $i^2 = -1$.
On dit que $z_M$ est l'\emph{affixe} du point $M$ et que $M$ est le \emph{point image} de $z_M$.

\begin{center}
\begin{tikzpicture}[scale=0.78,>=Stealth]
  \draw[gray!35,very thin,step=1] (-4,-4) grid (4,4);
  \draw[->,thick,bleu] (-4.5,0)--(4.5,0) node[right]{$\vec{u}$};
  \draw[->,thick,bleu] (0,-4.5)--(0,4.5) node[above]{$\vec{v}$};
  \node[below left,inner sep=2pt] at (0,0) {$O$};
  \foreach \x in {-4,...,-1} \node[below,font=\scriptsize,bleu] at (\x,0) {$\x$};
  \foreach \x in {1,...,4}   \node[below,font=\scriptsize,bleu] at (\x,0) {$\x$};
  \foreach \y in {-4,...,-1} \node[left,font=\scriptsize,bleu] at (0,\y) {$\y$};
  \foreach \y in {1,...,4}   \node[left,font=\scriptsize,bleu] at (0,\y) {$\y$};
  \fill[red] (3,2)   circle (2.4pt) node[above right,black,font=\small]{$A$};
  \fill[red] (-2,1)  circle (2.4pt) node[above left,black,font=\small]{$B$};
  \fill[red] (0,-3)  circle (2.4pt) node[right,black,font=\small]{$C$};
  \fill[red] (-1,-2) circle (2.4pt) node[left,black,font=\small]{$D$};
\end{tikzpicture}\\[2pt]
{\footnotesize Figure 1}
\end{center}

\begin{enumerate}[label=\textbf{\arabic*.}]
  \item Lire les coordonnées des points $A$, $B$, $C$ et $D$ de la figure 1.
  \item Écrire les affixes $z_A$, $z_B$, $z_C$ et $z_D$ sous la forme $a + ib$
        ($a$ et $b$ réels).
  \item Placer sur la figure les points $E$ et $F$ d'affixes respectives
        $z_E = 2 - 2i$ et $z_F = -3 + 3i$.
  \item On note $S$ le point d'affixe $z_S = z_A + z_B$.
    \begin{enumerate}
      \item Calculer $z_S$.
      \item Placer $S$, puis tracer les segments $[OA]$, $[OB]$ et $[OS]$.
      \item Comparer le vecteur $\overrightarrow{OS}$ avec $\overrightarrow{OA}+\overrightarrow{OB}$.
            Que constate-t-on ?
    \end{enumerate}
  \item Calculer $z_B - z_A$. Comparer avec les coordonnées du vecteur $\overrightarrow{AB}$.
  \item On appelle \emph{conjugué} de $z = a+ib$ le nombre $\bar{z} = a - ib$.
    \begin{enumerate}
      \item Donner $\bar{z_A}$ et $\bar{z_B}$.
      \item Placer les points $A'$ et $B'$, images de $\bar{z_A}$ et $\bar{z_B}$.
            Quelle transformation du plan relie $A$ et $A'$ ?
    \end{enumerate}
  \item Calculer le produit $z_A \times z_B$ en utilisant $i^2 = -1$ et en regroupant les termes.
\end{enumerate}

%=====================================================================
\section*{Partie B -- Module et argument : la forme trigonométrique (F.\,tr.)}

On considère les huit points de la figure 2.

\begin{center}
\begin{tikzpicture}[scale=1.25,>=Stealth]
  \draw[gray!35,very thin,step=1] (-2,-2) grid (2,2);
  \draw[->,thick,bleu] (-2.4,0)--(2.4,0) node[right]{$\vec{u}$};
  \draw[->,thick,bleu] (0,-2.4)--(0,2.4) node[above]{$\vec{v}$};
  \draw[thick,orange] (0,0) circle (1);
  \draw[thick,orange,dashed] (0,0) circle (1.414);
  \node[below left,inner sep=2pt] at (0,0) {$O$};
  \fill[red] (1,0)   circle (2pt) node[below right,black,font=\scriptsize]{$P$};
  \fill[red] (0,1)   circle (2pt) node[above right,black,font=\scriptsize]{$Q$};
  \fill[red] (-1,0)  circle (2pt) node[below left,black,font=\scriptsize]{$R$};
  \fill[red] (0,-1)  circle (2pt) node[below right,black,font=\scriptsize]{$S$};
  \fill[red] (1,1)   circle (2pt) node[above right,black,font=\scriptsize]{$T$};
  \fill[red] (-1,1)  circle (2pt) node[above left,black,font=\scriptsize]{$U$};
  \fill[red] (-1,-1) circle (2pt) node[below left,black,font=\scriptsize]{$V$};
  \fill[red] (1,-1)  circle (2pt) node[below right,black,font=\scriptsize]{$W$};
\end{tikzpicture}\\[2pt]
{\footnotesize Figure 2}
\end{center}

\begin{enumerate}[label=\textbf{\arabic*.}]
  \item Calculer les distances $OP$, $OQ$, $OR$, $OS$, puis $OT$, $OU$, $OV$, $OW$
        (valeurs exactes).
  \item Comparer ces résultats avec le nombre $\sqrt{a^2+b^2}$ où $z = a+ib$ est l'affixe du point.
        Ce nombre s'appelle le \textbf{module} de $z$ et se note $|z|$.
  \item Pour chaque point, déterminer une mesure de l'angle $\theta$ formé par l'axe des abscisses
        et la droite $(OM)$. Ce nombre s'appelle un \textbf{argument} de $z$.
  \item Compléter le tableau ci-dessous (le point $P$ est donné en exemple).
\end{enumerate}

\begin{center}
\small
\renewcommand{\arraystretch}{1.55}
\begin{tabular}{|c|c|c|c|>{\centering\arraybackslash}p{5.3cm}|}
\hline
\textbf{Point} & \textbf{Affixe } $z$ & \textbf{Module } $|z|$ & \textbf{Argument } $\theta$
& \textbf{Forme trigonométrique}\\
\hline
$P$ & $1$   & $1$ & $0$        & $1\,(\cos 0 + i\sin 0)$\\
\hline
$Q$ & $i$   &     &            & \\
\hline
$R$ & $-1$  &     &            & \\
\hline
$S$ & $-i$  &     &            & \\
\hline
$T$ & $1+i$ &     &            & \\
\hline
$U$ & $-1+i$&     &            & \\
\hline
$V$ & $-1-i$&     &            & \\
\hline
$W$ & $1-i$ &     &            & \\
\hline
\end{tabular}
\end{center}

\begin{enumerate}[label=\textbf{\arabic*.},start=5]
  \item Vérifier à la calculatrice que pour le point $T$ :
        \[ 1+i = \sqrt{2}\left(\cos\frac{\pi}{4} + i\sin\frac{\pi}{4}\right). \]
  \item \textbf{Bilan.} Compléter : tout nombre complexe non nul $z = a+ib$ peut s'écrire
        $z = r(\cos\theta + i\sin\theta)$ avec $r = \dots\dots$ et $\theta$ tel que
        $\cos\theta = \dots\dots$ et $\sin\theta = \dots\dots$
\end{enumerate}

%=====================================================================
\section*{Partie C -- Une transformation étonnante}

On reprend $z_A = 3+2i$ et $z_C = -3i$.

\begin{enumerate}[label=\textbf{\arabic*.}]
  \item Calculer $i\,z_A$ et $i\,z_C$ (on rappelle que $i^2=-1$).
  \item Placer les points $A_1$ et $C_1$, images respectives de $i\,z_A$ et $i\,z_C$.
  \item Comparer $|z_A|$ et $|i\,z_A|$, puis les arguments de $z_A$ et de $i\,z_A$.
  \item Recopier et compléter : « multiplier un nombre complexe par $i$, c'est faire tourner
        son point image de \dots\dots\dots\dots\dots\ autour de $O$, sans changer \dots\dots\dots\dots\dots. »
\end{enumerate}

\begin{center}
\begin{tikzpicture}[scale=0.78,>=Stealth]
  \draw[gray!35,very thin,step=1] (-3,-2) grid (3,3);
  \draw[->,thick,bleu] (-3.5,0)--(3.5,0);
  \draw[->,thick,bleu] (0,-2.4)--(0,3.5);
  \node[below left,inner sep=2pt] at (0,0) {$O$};
  \draw[dashed] (0,0)--(3,2);
  \draw[dashed] (0,0)--(-2,3);
  \draw[->,orange,thick] (33:1.0) arc (33:123:1.0);
  \node[orange,font=\scriptsize] at (0.35,1.35) {$90^\circ$};
  \fill[red] (3,2)  circle (2.4pt) node[below right,black,font=\small]{$A$};
  \fill[vert] (-2,3) circle (2.4pt) node[above left,black,font=\small]{$A_1$};
\end{tikzpicture}\\[2pt]
{\footnotesize Figure 3}
\end{center}

%=====================================================================
\section*{Partie D -- Application STI2D : l'impédance complexe}

En électricité, on représente une impédance par un nombre complexe
$Z = R + jX$ (on note $j$ au lieu de $i$ pour ne pas confondre avec l'intensité).
$R$ est la résistance et $X$ la réactance, exprimées en ohms ($\Omega$).

On considère un dipôle tel que $R = 50\ \Omega$ et $X = 50\sqrt{3}\ \Omega$.

\begin{enumerate}[label=\textbf{\arabic*.}]
  \item Calculer $|Z| = \sqrt{R^2 + X^2}$. C'est l'impédance du dipôle, en ohms.
  \item Calculer $\tan\varphi = \dfrac{X}{R}$, puis donner la valeur de $\varphi$ en radians
        et en degrés ($0 \leqslant \varphi \leqslant \frac{\pi}{2}$).
        Ce nombre $\varphi$ est le \emph{déphasage} du dipôle.
  \item Écrire $Z$ sous forme trigonométrique.
\end{enumerate}

%=====================================================================
\section*{Bilan de l'activité}

\begin{encadre}[À retenir]{vert}
Un nombre complexe s'écrit sous \textbf{forme algébrique} $z = a + ib$, où $a$ est la partie
\dots\dots\dots\dots\dots\dots{} et $b$ la partie \dots\dots\dots\dots\dots\dots{} .
Le point $M(a\,;b)$ est l'image de $z$ ; on dit que $z$ est l'\dots\dots\dots\dots{} de $M$.\\[4pt]
Le \textbf{module} $|z| = \sqrt{a^2+b^2}$ est la distance \dots\dots\dots\dots\dots\dots{} .\\
Un \textbf{argument} $\theta$ est une mesure de l'angle \dots\dots\dots\dots\dots\dots{} .\\[4pt]
On obtient la \textbf{forme trigonométrique} : \quad
$\boxed{z = r(\cos\theta + i\sin\theta)}$ \quad avec $r = |z|$.\\[4pt]
Multiplier par $i$, c'est effectuer une \dots\dots\dots\dots\dots\dots{} de centre $O$
et d'angle \dots\dots\dots\dots{} .
\end{encadre}



\end{document}